\section{A certified positive real period with exactly two vanishing factors}
\label{sec:real-pair-zero}
We construct an actual positive real period of the original geometric
coefficient family at which the symbol has order exactly two. The
point is defined by a unique root of the complete infinite series,
with explicit enclosing inequalities and a real two-variable proof.
The quartet, unit and conductor are those of PCL1--10 and WCF1--6
\cite{PCLOriginal,WCF}. The original four ES labels and their human
sources remain FC1--15a and \cite{ES488}. The geometric parameters
chosen here are not asserted to be arithmetic zeta zeros.

\subsection{The full original family and the exact real rectangle}
Retain
\[
 \delta=\frac14,\quad\gamma=3,\quad\beta=\frac{143}{8},
 \quad\eta=\frac{21025}{256},\qquad
 w=(\delta+i\gamma,\delta-i\gamma,-\delta+i\gamma,-\delta-i\gamma),
 \quad \rho_a=\frac12+w_a .\tag{RPZ1}
\]
Here $V_{ar}=w_a^r$, $0\le r\le3$, and the entire original period
matrix is
\[
 \mathcal R(z)=\sum_{n=0}^{\infty}R_nz^n,\qquad
 (R_n)_{ar}=\sum_{\substack{p,h\ge0\\2p+4h=5n+r+1-a}}
 \frac{(\beta/3)^p\eta^h(-5)^{p+h-n}(a/5)_{p+h-n}}{p!h!},
 \quad 1\le a\le4.\tag{RPZ2}
\]
The row index $a$ in this sum is distinct from the actual unit below.
Put $J_d=\operatorname{diag}(1,-1,-1,1)$ and $W(x)=I+ixJ_d$.
For real $x$ the original unit is $a=\alpha(1+ix)$, with its retained
nonzero real scale $\alpha$, and
$U=\alpha W(x)=\operatorname{diag}(a,\bar a,\bar a,a)$.
Its magnitude is $|\alpha|\sqrt{1+x^2}$. The common scalar cancels
with its inverse only in the original matrix conjugation.
For $x\in\mathbb C\setminus\{i,-i\}$ the same exact rational
matrices give the analytic extension
\[
 \begin{aligned}
 W(x)^{-1}&=\frac{I-ixJ_d}{1+x^2},\quad
 D_j=\operatorname{diag}(\zeta^{-j},\zeta^{-2j},\zeta^{-3j},\zeta^{-4j}),
 \quad\zeta=e^{2\pi i/5},\\
 M_j(z,x)&=W(x)^{-1}V\mathcal R(z)^{-1}D_j\mathcal R(z)V^{-1}W(x),\\
 f_j(\xi,z,x)&=Q_0\bigl(M_j(z,x)(e^{\rho_a\xi})_a\bigr),
 \quad Q_0(v)=v_1v_4-v_2v_3,\\
 d_j(z,x)&=f_j(0,z,x),\quad
 E_A(\xi,z,x)=\prod_{j=1}^4 f_j(\xi,z,x),\quad
 \mu_n(z,x)=\partial_\xi^n E_A(0,z,x).
 \end{aligned}\tag{RPZ3}
\]
The displayed extension follows by multiplication of $W$ and its
inverse; it is not restricted to UPP2's smaller disk. In particular
the real rectangle below satisfies $1+x^2>3$. The original PCL7--9
identity with the complete period conjugation, including the branch
powers and Gamma ray factors, is unchanged.

The following terminating decimals are exact rational numbers:
\[
 \begin{aligned}
 z_c&=0.02211947780866184513200791639903555836071244787217846841,\\
 x_c&=-1.466817440463494183975448871845763045856183562396293859,\\
 r_z&=r_x=10^{-15},\quad c_0=(z_c,x_c),\\
 X&=[z_c-r_z,z_c+r_z]\times[x_c-r_x,x_c+r_x]\subset\mathbb R^2 .
 \end{aligned}\tag{RPZ4}
\]
There is a unique $(z_*,x_*)\in X$ with $d_1(z_*,x_*)=0$.
It lies in the interior, and defines the exact unit
$a_*=\alpha(1+ix_*)$ and the positive original period $u_*=1/z_*$.
The proof of this existence and uniqueness assertion is given next.

\subsection{Full tails and derivatives used in the enclosing calculation}
The entrywise estimate FPZ5--6, proved directly from the coefficients
in RPZ2, is
\[
 |(R_n)_{ar}|\le C_R(2n+2)\frac{T_R^n}{n!},\qquad
 K=11/2,\quad C_R=4K^3,\quad T_R=6K^5/5 .\tag{RPZ5}
\]
For completeness its inequalities use $L=p+h-n\ge0$,
$(a/5)_L\le L!$,
$n!L!/(p!h!)=\binom{p+h}{p}/\binom{p+h}{n}\le2^{p+h}$,
$2(p+h)\le5n+3$, $2^{p+h}\le4\,6^n$, and at most $2n+2$
contributing pairs. The inequalities $K^2\ge5\beta/3$ and
$K^4\ge5\eta$ then give RPZ5 with every original sign retained
in RPZ2. This majorant proves convergence of the matrix and all
its derivative series on bounded sets.

Set $N=650$, let $\varrho_z$ be an outward upper bound for $|z|$ on the
real interval in RPZ4, and put $q=T_R\varrho_z$. The complete omitted
tails of the matrix and its first two derivatives are bounded by
\[
 \begin{aligned}
 \tau_0&=\frac{C_R(2N+4)q^{N+1}}{(N+1)!(1-2q/(N+2))},\\
 \tau_1&=\frac{C_R(2N+4)(N+1)T_Rq^N}{(N+1)!(1-3q/(N+1))},\\
 \tau_2&=\frac{C_R(2N+4)N(N+1)T_R^2q^{N-1}}
 {(N+1)!(1-q(N+3)/(N(N+2)))} .
 \end{aligned}\tag{RPZ6}
\]
These follow by summing the first omitted majorant with the
respective successive-term ratios bounded by
$2q/(N+2)$, $3q/(N+1)$ and $q(N+3)/(N(N+2))$.
All three bounds are less than one in the actual calculation.
Simultaneous Horner recurrence evaluates the three finite sums;
the respective tails are then added to each entry's real and
imaginary interval. The computed upper bounds are
\[
 \tau_0<2.099\,10^{-161},\quad
 \tau_1<9.485\,10^{-157},\quad
 \tau_2<1.349\,10^{-152}.\tag{RPZ7}
\]
No omitted coefficient is set to zero. A centered Taylor integral
further encloses each complete entry by
\[
 \begin{aligned}
 |\mathcal R_{ar}(z)-\mathcal R_{ar}(z_c)
              -\mathcal R'_{ar}(z_c)(z-z_c)|&\le L_{ar}|z-z_c|^2/2,\\
 |\mathcal R'_{ar}(z)-\mathcal R'_{ar}(z_c)|&\le L_{ar}|z-z_c|,
 \end{aligned}\tag{RPZ8}
\]
where $L_{ar}$ bounds the complete $|\mathcal R''_{ar}|$ on the
interval. Every integration segment stays in that interval.
Outward intervals enclosing the exact decimal center and radii
are used throughout, so rounding the center is also included.

Let $A=\mathcal R^{-1}$ and $C_j=AD_j\mathcal R$.
The derivative matrices are the literal products
\[
 \begin{aligned}
 (C_j)_z&=AD_j\mathcal R_z-A\mathcal R_zAD_j\mathcal R,\\
 (C_j)_{zz}&=AD_j\mathcal R_{zz}-2A\mathcal R_zAD_j\mathcal R_z
       +2A\mathcal R_zA\mathcal R_zAD_j\mathcal R
       -A\mathcal R_{zz}AD_j\mathcal R,\\
 (M_j)_{z^h}&=W^{-1}V(C_j)_{z^h}V^{-1}W\quad(h=0,1,2),\\
 (M_j)_x&=i[M_j,J_d]/(1+x^2),\qquad
 (M_j)_{zx}=i[(M_j)_z,J_d]/(1+x^2),\\
 (M_j)_{xx}&=-\{[[M_j,J_d],J_d]+2ix[M_j,J_d]\}/(1+x^2)^2.
 \end{aligned}\tag{RPZ9}
\]
Differentiating $A\mathcal R=I$ gives
$A_z=-A\mathcal R_zA$ and
$A_{zz}=2A\mathcal R_zA\mathcal R_zA-A\mathcal R_{zz}A$,
which proves the first two lines with their stated matrix orders.
For the phase derivative use
$W^{-1}W_x=(xI+iJ_d)/(1+x^2)$; its scalar term commutes with
$M_j$. Differentiation gives the remaining lines. With the original
bilinear polar form
$\mathfrak b(v,w)=v_1w_4+v_4w_1-v_2w_3-v_3w_2$,
$y=M_j\boldsymbol1$ and $h=M_j(\rho_a)_a$, the scalar derivatives
are $d_j=\mathfrak b(y,y)/2$,
$d_{j,z}=\mathfrak b(y,y_z)$,
$d_{j,x}=\mathfrak b(y,y_x)$ and
$f_{j,\xi}=\mathfrak b(y,h)$. The second derivatives used to
enclose them follow by differentiating these products, exactly
UPP6. In particular no Hermitian form replaces this bilinear form.

\subsection{A completely specified real contraction}
Define the real map $F(z,x)=(\Re d_1(z,x),\Im d_1(z,x))^{\mathsf T}$.
Use the fixed exact dyadic matrix
\[
 P=2^{-64}\begin{pmatrix}
 3218485790678941&-57245565561083039\\
 -5877924768298298697&-18277625187819728703
 \end{pmatrix},\qquad T(y)=y-PF(y).\tag{RPZ10}
\]
The enclosing calculation is supplied as
\nolinkurl{certify_real_pair.py}.
Its receipt is \nolinkurl{REAL_PAIR_CERTIFICATE.json}.
It uses Python-FLINT 0.9.0,
768-bit inclusion arithmetic and one worker. The arithmetic source
is Fredrik Johansson's original Arb paper \cite{human:arb2016},
Introduction and the sections on radii and precision. The present
tail, derivative and fixed-point proof is given here, rather than
attributed to that arithmetic reference.

The complete original matrix is invertible throughout $X$.
Specifically, with $A_c$ the dyadic midpoint matrix of the enclosing
inverse at $z_c$, the full matrix enclosure gives
$\|I-A_c\mathcal R(z)\|_\infty<1.249\,10^{-11}<1$.
The convergent geometric matrix series in $I-A_c\mathcal R(z)$
inverts $A_c\mathcal R(z)$, so its two square factors, including
$\mathcal R(z)$, are invertible. The inverse evaluation is enclosing.
The determinant of RPZ10's exact matrix satisfies
\[
 -0.001161715807429<\det P<-0.001161715807427 .\tag{RPZ11}
\]
For reference, the full center Jacobian is enclosed about
$\left(\begin{smallmatrix}852.9039895042&-2.671297323736\\
-274.2864805121&-0.150186872902\end{smallmatrix}\right)$;
all intervals, rather than this shortened display, enter the proof.

Let $\boldsymbol J$ enclose $DF(y)$ for all $y\in X$, let
$\boldsymbol E=I-P\boldsymbol J$, and let $\boldsymbol\eta$
enclose $PF(c_0)$. For component indices $i,j\in\{z,x\}$ put
\[
 q_i=\frac{\sum_j\sup|\boldsymbol E_{ij}|r_j}{r_i},\qquad
 b_i=\frac{\sup|\boldsymbol\eta_i|
                    +\sum_j\sup|\boldsymbol E_{ij}|r_j}{r_i}.
 \tag{RPZ12}
\]
The complete enclosures verify the strict rational inequalities
\[
 q_z,b_z<4.578\,10^{-9},\qquad
 q_x,b_x<1.829\,10^{-6}<1.\tag{RPZ13}
\]
All operations defining the enclosures, the full defect matrix,
center residual, exact dyadic $P$, and tail bounds are retained in
the executable and receipt. They are sufficient to replay each
displayed strict inequality without any numerical root assumption.

Here is the existence proof. Equip $\mathbb R^2$ with
$\|v\|_r=\max_i|v_i|/r_i$. Integration of $DT=I-PDF$ along the
segment in the convex rectangle $X$ gives
$\|T(y)-T(y')\|_r\le q\|y-y'\|_r$, where
$q=\max_iq_i<1$. Also
$T(c_0)-c_0=-PF(c_0)$, so the same integration gives
$|T(y)_i-(c_0)_i|\le b_ir_i<r_i$.
Thus $T$ maps $X$ into its interior and contracts it.
Start $y_0=c_0$ and set $y_{n+1}=T(y_n)$. Induction keeps every
iterate in $X$, and
\[
 \|y_m-y_n\|_r\le
 \frac{q^n}{1-q}\|y_1-y_0\|_r\quad(m>n).\tag{RPZ14}
\]
This follows by summing the geometric bounds on successive
differences. Hence the iterates form a Cauchy sequence. Completeness
of $\mathbb R^2$ and closedness of $X$ give a limit $y_*\in X$.
Continuity gives $T(y_*)=y_*$, and RPZ11 gives $F(y_*)=0$.
Any two fixed points have distance at most $q$ times their distance,
so they coincide. Every zero of $F$ is a fixed point. Finally
$y_*=T(y_*)$ lies in the interior by the strict self-map bound.
This proves every existence and uniqueness statement after RPZ4.

\subsection{The exact pair, moments and nonzero inverse coefficient}
On the entire real rectangle the factor enclosures include
\[
 \begin{aligned}
 d_{1,z}&\in[852.90399\pm1.90\,10^{-6}]
                   +i[-274.28648\pm1.91\,10^{-6}],\\
 f_{1,\xi}&\in[17.1783460\pm3.61\,10^{-8}]
                   +i[-6.8131052\pm4.44\,10^{-8}],\\
 d_2&\in[0.05410876\pm1.84\,10^{-9}]
                   +i[-0.31535702\pm5.02\,10^{-9}].
 \end{aligned}\tag{RPZ15}
\]
They exclude zero. Let $J_c$ exchange the first two and the last
two primary coordinates; it is distinct from $J_d$ in RPZ3.
The real coefficient matrix in RPZ2, $\bar V=J_cV$,
$\bar U=J_cUJ_c$, and $Q_0(J_cv)=-Q_0(v)$ give
\[
 f_{5-j}(\xi,z,x_*)=-\overline{f_j(\bar\xi,\bar z,x_*)},\qquad
 E_A=G G^\#,\quad G=f_1f_2,\quad
 G^\#(\xi,z)=\overline{G(\bar\xi,\bar z)}.\tag{RPZ16}
\]
This is the earlier OFM8--11 conjugation and positivity identity
\cite{PCLOriginal,Cumulative}, with its full two-variable use
derived in RPCJ1--4 below. It fixes the original actual unit.
At the certified point exactly $d_1=d_4=0$ and $d_2,d_3\ne0$.
Writing $A=G_z(0,z_*,x_*)$ and $B=G_\xi(0,z_*,x_*)$, RPZ15
gives $A,B\ne0$. Consequently the full analytic products give
\[
 \begin{aligned}
 \mu_0(z,x_*)&=|A|^2(z-z_*)^2+O((z-z_*)^3),\\
 \mu_0(z_*,x_*)&=\mu_1(z_*,x_*)=0,\quad
 \mu_2(z_*,x_*)=2|B|^2>0 .
 \end{aligned}\tag{RPZ17}
\]
Thus the period zero has multiplicity exactly two and the original
symbol has order exactly two, in the original coordinates.
Every nontrivial orbit quadric has nonzero restriction to the
exponential curve: for $j=1,4$ its first derivative is nonzero,
and for $j=2,3$ its value is nonzero. Since $Q_0(t(\xi))=0$
identically, none is associated to $Q_0$. If two orbit quadrics
were associated, applying the inverse cyclic action would associate
a nontrivial member to $Q_0$. Therefore all five orbit quadrics
are distinct. This does not assert pairwise coprimality after every
unrelated section.

The full original derivatives, not freely chosen jets, satisfy
\[
 \begin{aligned}
 A&\in[-40.34859\pm1.87\,10^{-6}]
                 +i[-283.81057\pm6.15\,10^{-6}],\\
 B&\in[-1.2190616\pm5.30\,10^{-8}]
                 +i[-5.7859608\pm6.51\,10^{-8}],\\
 c=|A|^2&\in[82176.45\pm0.00525],\qquad
 b=|B|^2\in[34.963453\pm4.49\,10^{-7}],\\
 R=2\Re(A\bar B)&\in[3382.6084\pm0.0000642],\qquad
 2cb-R^2\in[-5695695\pm0.547].
 \end{aligned}\tag{RPZ18}
\]
In particular $R(2cb-R^2)\ne0$. The complete original moment
enclosures through order five are supplied by RPE2--3 below; their
exact finite derivative formula retains all four factors. At this
actual point the generic inverse coefficient of RPCJ15--16 is
nonzero. Hence the fixed old degree-three conductor has inverse
singular exponents $(5,3,0,0)$ and inverse exterior exponents
$(5,8,8,8)$, with every original Gamma constant calculated in
RPCJ20--22 and RPE5--9. The exterior rank-one leading coefficient
before its original factor $\rho_3/\rho_0$ is
\[
 \frac{|R(2cb-R^2)|}{c^4}
       \in[4.224828\,10^{-10}\pm7.04\,10^{-17}],\qquad
 u_*\in[45.20902385898\pm3.31\,10^{-12}].\tag{RPZ19}
\]
The source order change, four-label return, and both lost directions
are calculated below on their full original spaces. This positive
real geometric period is an actual obstruction to a finite uniform
inverse bound on the whole geometric family. Locating the arithmetic
quartet and its restricted period/unit locus requires their exact
arithmetic map; the present construction makes no identification of
that locus with this explicitly certified member.
